\section{因果术语与统计术语之辨}

本节定义贯穿全书的若干基本术语与概念。这些定义可能与标准文献中的定义不尽一致，因此，凡对这些术语的解释存有疑问，都应以本节为准。

\emph{概率参数}是任何以联合概率函数定义\footnote{若量 $Q$ 可由 $C$ 类中任一对象的描述唯一算出（即 $Q$ 是由 $C$ 到 $Q$ 之论域的一个函数映射），则称 $Q$ 以 $C$ 类对象定义。}的量。例子见 \S 1.1 与 \S 1.2 中所定义的各个量。

\emph{统计参数}是任何以\emph{被观测}变量的联合概率分布定义的量，它对未观测变量存在与否不作任何假设。

\noindent 例如：条件期望 $E(Y\mid x)$；回归系数 $r_{YX}$；密度函数在 $y=0,x=1$ 处的值。

\emph{因果参数}是任何以因果模型（如 \textup{(1.40)}）定义、且不是统计参数的量。

\noindent 例如：\textup{(1.41)} 中的系数 $\beta_{ik}$；对某个 $u$，$X_9$ 是否对 $X_3$ 有影响；在干预 $\dox{X=0}$ 之下 $Y$ 的期望值；变量 $X_7$ 的父节点数目。

\begin{remarkenv}
把未测量变量排除在统计参数的定义之外，是为了防止人们把因果假设藏在潜变量的外衣之下。若允许这类构造，任何量都可以被说成是统计的，于是“仅凭统计数据即可估计的量”与“需要数据之外的额外假设才能确定的量”之间那重要区别，就被模糊了。
\end{remarkenv}

\emph{统计假设}是对被观测变量之联合分布的任何约束；例如，$f$ 为多元正态，或 $P$ 相对于给定 DAG $D$ 为马尔可夫的。

\emph{因果假设}是对因果模型的任何约束，且该约束不能通过施加统计假设来实现；例如，$f_i$ 是线性的，$U_i$ 与 $U_j$（未被观测）不相关，或者 $x_3$ 不出现在 $f_4(\pa_4,u_4)$ 中。因果假设可能有、也可能没有统计含义。有统计含义时，我们称该假设是“可检验的”或“可否证的”。因果假设常常（虽非总是）可以由实验研究予以否证，此时我们称它们是“实验可检验的”。例如，“$X$ 对图 1.6 模型 2 中的 $E(Y)$ 没有影响”这一假设是经验可检验的；而“$X$ 可能治愈总体中某位特定受试者”这一假设则不是。

\begin{remarkenv}
因果参数与统计参数之别是明确而根本的——二者不可混同。除非援引因果假设，否则无法从统计参数中识别出因果参数。这些假设的表述与简化，将占据本书的主要篇幅。
\end{remarkenv}

\begin{remarkenv}
变量之间的时间先后可能提供关于因果关系（之缺失）的某些信息——较晚的事件不可能是较早事件的原因。因此，诸如 $P(y_t\mid y_{t-1},x_t),\ t=1,\dots$ 这样按时间标号的分布（它们在经济分析中被例行使用），或许可视为统计模型与因果模型之间的边缘情形。但我们仍将这类模型归入统计模型，因为在我们承诺不对未测量变量之有无作任何假设的前提下，绝大多数与政策相关的问题都无法从这类分布中识别出来。因此，“格兰杰因果”（Granger 1969）与“强外生性”（Engle et al. 1983）这类计量经济学概念，将被归入统计概念而非因果概念。\footnote{在把概率序列 $P(y_t\mid y_{t-1},x_t),\ t=1,\dots$ 称为“数据生成模型”时也须谨慎（如 Davidson and MacKinnon 1993, p.~53; Hendry 1995）。这类序列本质上是统计性的；除非援引第二章所发展起来的那类因果假设（见定义 2.4.1 与 2.7.4），否则它们不能用于政策评价任务。}
\end{remarkenv}

\begin{remarkenv}
“理论的”与“结构的”二词常与“因果的”互换使用；我们将使用后两个词，同时记住有些结构模型未必是因果的（见 \S 7.2.5）。
\end{remarkenv}

\subsection*{因果概念与统计概念}

因果参数与统计参数之间的分界线同样延伸到一般概念，并将由术语上的区分加以支撑。统计概念的例子有：相关、回归、条件独立、关联、似然、可压缩性、风险比、优势比、倾向评分、格兰杰因果，等等。因果概念的例子有：随机化、影响、效应、混杂、外生性、可忽略性、扰动（如 \textup{(1.40)}）、伪相关、路径系数、工具变量、干预、解释，等等。划这条界线的目的，不是要把因果概念排除在统计分析的范围之外，而是鼓励研究者用恰当的工具去处理非统计的概念。

有些读者也许会感到意外：随机化、混杂、伪相关、效应这类教科书概念竟然是\emph{非统计}的。另一些读者也许会对下述想法感到震惊：外生性、混杂、反事实这类颇有争议的概念，竟可以用因果模型来定义。本书正是为这些读者而写，接下来的篇章将表明，刚刚作出的因果概念与统计概念之别，对于澄清二者都是必不可少的。

\subsection*{因果分析的两道心理障碍}

统计概念与因果概念之间的鲜明区分，可以转述为一条有用的原则：\textbf{每一个因果断言背后，都必定隐藏着某个无法从联合分布中识别、因而无法在观测研究中检验的因果假设。}这类假设通常由人依据专家判断提供。于是，人类组织并交流经验知识的方式，便成为本研究不可分割的一部分，因为它决定了要求专家表述的那些判断是否可靠。

这一因果—统计之分的另一层含义是：任何因果分析的数学进路都必须引入新的记号。概率演算的词汇——连同其强大的期望、条件化与边缘化算子——严格地以分布函数来定义，因而不足以表达因果假设或因果断言。为作说明，概率演算的句法不允许我们表达“症状不会导致疾病”这一简单事实，更不用说从这类事实中得出数学结论了。我们所能说的只是两个事件相依——意思是若发现其一，便可预期遇到另一个；但我们无法区分由条件概率 $P(\text{疾病}\mid\text{症状})$ 所量化的统计相依，与因果相依，因为后者在标准概率演算中没有表达式。

前述两项要求——(1) 以未经检验的判断性假设开始因果分析；(2) 扩展概率演算的句法——构成了传统统计学训练有素的专业人士接受因果分析的两大障碍（Pearl 2003c；另见 \S 11.1.1 与 \S 11.6.4）。本书借助一套基于图形进路与代数进路共生的有效而友好的记号系统，帮助读者跨越这两道障碍。
